\begin{pastexam}[label={px:ncku-114-isec-ds-03}]{成大智慧資安學程 114 學年度 單選題第 3 題}
（單選；每題 5 分，答錯倒扣 2 分）

Given a directed graph in the following, please identify which edge would become a cross edge if we
traverse the graph by the breadth first search algorithm from node C. Assume that the nodes are stored
in the adjacent list according to alphabetic order. Please specify which edge forms a cross edge in the
resulting breadth-first spanning tree.

\begin{center}
\begin{tikzpicture}[v/.style={draw,circle,minimum size=7mm,inner sep=0pt},>={Stealth[length=2.2mm]}]
  \node[v] (A) at (0,0) {A};
  \node[v] (B) at (1.7,0.9) {B};
  \node[v] (C) at (2.1,-0.8) {C};
  \node[v] (D) at (4.1,0) {D};
  \node[v] (E) at (6.0,0) {E};
  \node[v] (F) at (6.0,1.5) {F};
  \draw[->] (A) -- (B);
  \draw[->] ([yshift=1.5mm]C.west) -- ([yshift=-0.5mm]A.east);
  \draw[->] (A.south east) -- ([yshift=-1.5mm]C.west);
  \draw[->] ([xshift=-1mm]B.south) -- ([xshift=-2.5mm]C.north);
  \draw[->] ([xshift=1mm]C.north) -- ([xshift=2.5mm]B.south);
  \draw[->] (B.east) -- (D.north west);
  \draw[->] ([yshift=2.5mm]D.north west) -- ([yshift=2.5mm]B.east);
  \draw[->] (C) -- (D);
  \draw[->] (D) -- (F);
  \draw[->] (D) -- (E);
  \draw[->] (F) -- (E);
\end{tikzpicture}
\end{center}
a)~(F, E) \qquad b)~(B, D) \qquad c)~(D, B) \qquad d)~(A, B) \qquad e)~All of the above.
\end{pastexam}

\begin{pastsoln}
答案是 e) All of the above。

圖的邊（照原圖的箭頭）：$A \to B$、$A \to C$、$B \to C$、$B \to D$、$C \to A$、$C \to B$、$C \to D$、
$D \to B$、$D \to E$、$D \to F$、$F \to E$。 鄰接串列照字母順序：
$A: B, C$；$B: C, D$；$C: A, B, D$；$D: B, E, F$；$E$：無；$F: E$。

\textbf{BFS 從 C 開始}（資結書第~21.3~節），$d$ 是層數：
\begin{center}
\small
\begin{tabular}{lll}
\toprule
取出 & 檢查的邊 & 佇列 \\
\midrule
（開始） & & C \\
C（$d = 0$） & $C \to A$、$C \to B$、$C \to D$ 都是新點，成為樹邊 & A, B, D \\
A（$d = 1$） & $A \to B$、$A \to C$ 都已發現 & B, D \\
B（$d = 1$） & $B \to C$、$B \to D$ 都已發現 & D \\
D（$d = 1$） & $D \to B$ 已發現；$D \to E$、$D \to F$ 是新點，成為樹邊 & E, F \\
E（$d = 2$） & 沒有出邊 & F \\
F（$d = 2$） & $F \to E$ 已發現 & \\
\bottomrule
\end{tabular}
\end{center}
廣度優先生成樹：C 的子節點是 A、B、D，D 的子節點是 E、F。

\textbf{分類非樹邊}（資結書第~21.5~節）：$(u, v)$ 不是樹邊時，$v$ 是 $u$ 的祖先就是後向邊，$v$ 是 $u$ 的
子孫就是前向邊（BFS 不會出現），兩者都不是就是交叉邊。
\begin{itemize}
\item $(A, C)$、$(B, C)$：C 是根，是 A、B 的祖先，後向邊。
\item $(A, B)$：A、B 都是 C 的子節點，互不為祖先，\textbf{交叉邊}。
\item $(B, D)$、$(D, B)$：B、D 是兄弟，\textbf{交叉邊}。
\item $(F, E)$：E、F 都是 D 的子節點，\textbf{交叉邊}。
\end{itemize}
選項 a) 到 d) 的四條邊都是交叉邊，答案 e)。
\end{pastsoln}
